\section{The pole-free closure is not a convolution submodule}
\label{sec:pole-free-module-obstruction}

The correction in MA-MO-011 gives an endomorphism on a genuine
Hilbert space, but the corrected domain is not preserved by the
translation representation. The failure holds for every nonzero
compact smooth convolution, including tests vanishing at logarithmic
zero. This supplies a precise obstruction to an equivariant
retraction onto that domain, without ruling out a different transport.

Keep \(J,K,\Lambda,v,\phi\) and their exact inherited norms
from MA-MO-011. Thus its fixed-gap smooth or generic-width Gevrey
regime is assumed for the cutoff and the two moment functionals.
In the present result \(g\in C_c^\infty(\mathbb R)\) may be
arbitrary in either regime: only boundedness of convolution is
needed, not the trace-class assertion for the raw smoothed defect.
Let
\[
 (U_\tau y)(x)=y(x-\tau),\qquad
 C_gy=\int g(z)U_zy\,dz.
 \tag{MA-MO.200}
\]

\paragraph{MA-MO-013 (no equivariant retraction onto the pole-free closure).}
Both \(U_\tau\) and \(C_g\) are bounded on the inherited
sum Hilbert space and preserve its closed theta ideal \(J\).
Nevertheless
\[
 \begin{aligned}
 U_\tau(K)&\not\subseteq K&&\text{for every }\tau\ne0,\\
 \Lambda C_g:K&\longrightarrow\mathbb C^2
     \text{ is onto}&&\text{for every }g\ne0.
 \end{aligned}
 \tag{MA-MO.201}
\]
The second assertion holds already on genuine compact additive
theta inputs in \(\mathscr U_{\rm pf}\). It is therefore
stronger than failure of invariance in the Hilbert closure.

No bounded retraction \(P:J\to J\) with
\(P^2=P\) and \(\operatorname{ran}P=K\) can commute with
any nontrivial \(U_\tau\), or with any nonzero \(C_g\)
of this compact smooth class. In particular the orthogonal
projection (MA-MO.140) is not a convolution-module retraction.
In fact the only linear subspace of \(K\) invariant under
all \(U_\tau\), \(\tau\in\mathbb R\), is \(\{0\}\).
Consequently any intertwiner from a representation of this
translation group into \(J\) whose range is contained in
\(K\) is zero.

\begin{proof}
\emph{Actual bounded representation and smooth source.}
For the weight \(w\) of (MA-MO.137), the ratio
\(w(x+\tau)/w(x)\) is bounded uniformly in \(x\) for
each fixed \(\tau\), locally uniformly in \(\tau\).
Indeed the exponential part is bounded by
\(e^{2a|\tau|}\), and
\(1+(x+\tau)^2\le2(1+x^2)(1+\tau^2)\)
bounds the polynomial factor. Thus translation and its
inverse are bounded, and integration over the compact
support of \(g\), followed by the triangle inequality
in the Hilbert norm, proves boundedness of \(C_g\).

On the even additive Schwartz source these operations are
\[
 \begin{aligned}
 \rho_\tau f(t)&=e^{-\tau/2}f(e^{-\tau}t),\\
 \mathcal D_g f(t)&=\int g(z)e^{-z/2}f(e^{-z}t)\,dz,\\
 U_\tau\Theta f&=\Theta\rho_\tau f,\qquad
 C_g\Theta f=\Theta\mathcal D_gf.
 \end{aligned}
 \tag{MA-MO.202}
\]
The sums and integrals interchange absolutely, with all
derivatives on compact norm intervals. These formulas also
follow directly from the actual theta sum in MA-MO.83.
Dilation and its compact integral preserve even Schwartz
space, the value zero at the additive origin and the zero
integral. Hence they preserve \(\mathscr U_0\), and
boundedness gives preservation of \(J\) by taking closure.

\emph{A compact smooth convolution cannot reproduce the cutoff.}
Use the actual functions \(W_g\) of (MA-MO.149). We claim
\[
        W_g\notin\mathbb C v
                      \quad(g\in C_c^\infty,\ g\ne0).
 \tag{MA-MO.203}
\]
If \(W_g=c v\), differentiate at \(t=e^x\) and take
the Fourier transform in \(x\), just as in (MA-MO.150).
It gives
\[
       \bigl(\mathcal Mg(1/2-it)-c\bigr)P(t)=0,
               \qquad P(t)=\widehat{\phi'}_+(t),\quad P(0)=-1.
 \tag{MA-MO.204}
\]
The entire function \(\mathcal Mg(1/2-it)\) therefore
equals \(c\) on an interval, hence everywhere. Along
the real \(t\) axis it is a Schwartz Fourier transform,
so it tends to zero at infinity; consequently \(c=0\).
Fourier injectivity then gives \(g=0\), a contradiction.
This proves (MA-MO.203), including the possibility \(c=0\),
and applies as well to \(g^\vee\).

\emph{Both outgoing convolution moments have full rank.}
Let \(W=W_g\) and \(U=W_{g^\vee}\). The bilinear
scaling and Fourier identities proved in MA-MO-011 give,
without the source operator \(L\),
\[
 \lambda_0(C_g\Theta f)=\int fW,\qquad
 \lambda_1(C_g\Theta f)=\int f\mathfrak FU
                        \quad(f\in\mathscr U_{\rm pf}).
 \tag{MA-MO.205}
\]
All functions here are ordinary smooth functions and all
pairings are absolutely convergent. We prove independence
of the five even functions
\[
                     1,\quad v,\quad\mathfrak Fv,
                              \quad W,\quad\mathfrak FU
                  \quad\text{on }\mathbb R\setminus\{0\}.
 \tag{MA-MO.206}
\]
Suppose \(A W+B\mathfrak FU=c_0+c_1v+c_2\mathfrak Fv\).
All terms other than \(c_0\) are Schwartz, so decay gives
\(c_0=0\). Outside the compact supports of \(W\) and
\(v\), the identity becomes
\(B\mathfrak FU=c_2\mathfrak Fv\). Both transforms
are entire, because \(U\) and \(v\) are compact smooth.
It follows by the identity theorem and Fourier inversion
that \(BU=c_2v\). If \(B\ne0\), this contradicts
(MA-MO.203) for \(g^\vee\). Thus \(B=0\) and
\(c_2=0\). The remaining identity \(AW=c_1v\)
contradicts (MA-MO.203) if \(A\ne0\). Hence \(A=0\)
and \(c_1=0\), proving the claimed independence.

As in MA-MO.97, choose five positive evaluation points
with an invertible matrix for these five functions,
replace each point and its negative by an even bump
of mass one supported away from zero, and take the
bumps narrow enough that their moment matrix remains
invertible. Given arbitrary \((z_0,z_1)\in\mathbb C^2\),
solve the resulting system with right side
\((0,0,0,z_0,z_1)^{\mathsf T}\). Its actual compact
even solution \(f\) has
\[
        f(0)=\int f=\ell_0(f)=\ell_1(f)=0,
                \qquad \int fW=z_0,\quad\int f\mathfrak FU=z_1.
 \tag{MA-MO.207}
\]
Equations (MA-MO.139) and (MA-MO.205) prove the second
assertion of (MA-MO.201), on the claimed smooth domain.
The invertible matrix exists because otherwise the span
of all evaluation vectors would have a nonzero annihilator,
contradicting (MA-MO.206); its bump approximation follows
from continuity of these functions. Thus the construction
does not rely on a distributional input.

\emph{Every nontrivial translation also fails.}
The first moment of \(U_\tau\Theta f\) pairs \(f\)
with \(v_\tau(t)=e^{\tau/2}v(e^\tau t)\).
If \(v_\tau\) belonged to
\(\operatorname{span}\{1,v,\mathfrak Fv\}\) away zero,
the exterior identity and decay would remove the constant
term, and entire continuation would remove the
\(\mathfrak Fv\) term. Thus \(v_\tau=c v\).
Near zero both undilated cutoff values are one, so
\(c=e^{\tau/2}\). This would imply
\(v(e^\tau t)=v(t)\) everywhere. Iterating the dilation
or its inverse takes a fixed positive argument outside
the compact support, contradicting \(v=1\) near zero.
Therefore the four functions \(1,v,\mathfrak Fv,v_\tau\)
are independent for \(\tau\ne0\). The same four-bump
construction gives \(f\in\mathscr U_{\rm pf}\) with
\(\lambda_0(U_\tau\Theta f)=1\). This proves the first
assertion of (MA-MO.201).

\emph{The retraction obstruction.}
If \(P\) is such a retraction and commutes with a
bounded operator \(T\) on \(J\), then for every
\(y\in K\) one has \(Ty=TPy=PTy\in K\).
Apply this to either class of operators in (MA-MO.201)
to obtain a contradiction. This proves the stated
impossibility for every proposed bounded retraction,
not only for the orthogonal one.

\emph{No nonzero invariant subspace remains inside the kernel.}
There is also a direct stronger argument, using the actual
functional constructed in MA-MO-011. Let \(y\in H_\Sigma\),
put \(u=V_ay\in L^2\), and retain \(\beta_a\) of
(MA-MO.143). Fourier scaling gives
\[
 \lambda_0(U_\tau y)
   =e^{a\tau}\int_{\mathbb R}u(t)\beta_a(t)e^{it\tau}\,dt.
 \tag{MA-MO.208}
\]
The product is integrable by Cauchy--Schwarz, since
\(\beta_a\in L^2\). Moreover \(\beta_a\ne0\) almost
everywhere. Indeed \(j_a\) is finite and nonzero on the
chosen line in both stated regimes. By (MA-MO.144), the
remaining factor is the nonzero entire transform of
\(e^{cx}\phi'(x)\), with \(c=1/2-a\), divided by a nonvanishing linear
factor on the real axis. The numerator is not identically
zero by Fourier uniqueness, and its real zeros are discrete.

If every translate of \(y\) belongs to \(K\), the left
side of (MA-MO.208) vanishes for every real \(\tau\).
Uniqueness of the Fourier transform of an \(L^1\) function
gives \(u\beta_a=0\) almost everywhere, hence \(u=0\)
and \(y=0\). Every element of a translation-invariant
linear subspace of \(K\) has that property, proving the
stronger assertion without a closedness assumption on
that subspace. If an intertwiner has range in \(K\),
the orbit of each of its values is also in \(K\), so
each value is zero. No claim about a different target
space or an unequivariant map is implied.
\end{proof}

\paragraph{Scope for an intertwiner and a trace.}
The original representation on \(J\) is well defined,
and MA-MO-011 still gives legitimate corrected endomorphisms
and determinants on \(K\). What fails is using a retraction
onto \(K\) as a translation- or convolution-equivariant
domain repair. In particular the rank-two correction is
not an intertwining argument for the original source
representation. This holds also for every nonzero compact
test with \(g(0)=0\); that pointwise condition does not
remove either independent outgoing moment.

The result concerns this precisely defined rational
pole-free closure in the two analytic regimes already
stated. It neither rules out all global-to-local transports
nor changes the exact arithmetic deficiency condition in
MA-FT-002. In the inherited cote~1 framework, scans~079--080
and~111--112 through BR-TR-01/02, existence of an ordinary
nuclear endomorphism trace is separate from compatibility
with a representation or its convolution algebra. The
finite-length \(K_0\) bookkeeping does not supply that
missing equivariance.
